\documentclass[10pt,a4paper]{amsart}
\usepackage[margin=19mm]{geometry}
\usepackage{amsmath,amssymb,mathtools}
\usepackage[scr=rsfs,cal=euler]{mathalfa}
\usepackage{xfrac}
\usepackage[unicode,hidelinks,bookmarks=false]{hyperref}
\newcommand{\ie}{i.e.\ }
\newcommand{\cf}{cf.\ }
\newcommand{\resp}{respectively\ }
\newcommand{\Subsection}{\subsection}
\providecommand{\nobreakdash}{\nobreak\hbox{-}}
\setlength{\emergencystretch}{2em}
\multlinegap0pt
\usepackage{bm}
\usepackage{bbm}
\usepackage{dsfont}
\usepackage{xspace}
\usepackage{cancel}
\usepackage{mathtools}
\usepackage[shortlabels]{enumitem}
\usepackage{amsthm}
\makeatletter
\@ifundefined{theorem}{\newtheorem{theorem}{Theorem}}{}
\@ifundefined{lemma}{\newtheorem{lemma}[theorem]{Lemma}}{}
\makeatother
\makeatletter
\newcommand{\N}{\mathbb{N}}
\newcommand*\diff{\mathop{}\!\mathrm{d}}
\newcommand{\e}{\mathrm{e}}
\newcommand{\lra}{\longrightarrow}
\expandafter\ifx\csname myproof\endcsname\relax
\newenvironment{myproof}[2] {\paragraph{\textsc{Proof of {#1}}.}}{\hfill$\square$\newline}
\fi
\makeatother
\title[A generalization of Gr\"unbaum's inequality in RCD$(0,N)$-sp]{A generalization of Gr\"unbaum's inequality in RCD$(0,N)$-spaces}
\author{Victor-Emmanuel Brunel, Shin-ichi Ohta, Jordan Serres}
\date{Source: arXiv:2408.15030v1 (CC BY 4.0)}
\begin{document}
\maketitle

\noindent\emph{Source and licence.} A generalization of Gr\"unbaum's inequality in RCD$(0,N)$-spaces. Version arXiv:2408.15030v1 (CC BY 4.0), distributed under the Creative Commons Attribution 4.0 International licence, as recorded in the arXiv licence metadata for this exact version and shown on the abstract page https://arxiv.org/abs/2408.15030v1. Full rights notice: licenses/arxiv-2408.15030-CC-BY-4.0.txt.\par\smallskip

\noindent\textit{Editorial scope.} Counted unit: the Lemma whose source label is \texttt{lm:=\_inf} in the article (source0.tex lines 999--1019, source label \texttt{lm:=\_inf}), delivered together with the article's own complete original proof of it (source0.tex lines 1021--1089, one \texttt{proof} environment of 69 lines). No mathematics is written, repaired or re-derived: statement and proof are the article's own text line for line. Required context retained uncounted: lm:1d\_neg (lines 915--932); eq:1d\_neg (lines 921--924); eq:1d\_inf (lines 927--930); eq:H\_neg (lines 966--969); eq:b\_neg (lines 972--977). Every definition, display and label of those lines is retained unchanged. Omitted from this package and not redistributed: 1--914, 925--926, 931--965, 970--971, 978--998, 1020--1020, 1090--1964. Nothing inside a retained excerpt is omitted or altered: every retained body is delivered exactly as the article's lines are written, so no omission receipt of any kind accompanies this package. Citations: the retained excerpts carry no citation command at all, so no bibliography entry of the article is transcribed and no reference is redistributed. Reference adaptation: none was needed and none was made; every reference inside the delivered unit resolves inside it, so no unresolved reference or citation is produced. Printed slips kept literal: no printed word, symbol, hypothesis or number of the counted statement or of its proof was altered. Layout and class: the article's own class is \texttt{imsart} with packages including latexsym, amssymb, amsmath, eucal, epic, eepic, color, enumerate, comment, hyperref, amsthm, cite, inputenc, dsfont; the wrapper is the lane's pinned \texttt{amsart} editorial wrapper at 10pt on A4, which restates the theorem-like environments the retained text needs and adds the article's own macro definitions that the retained text needs (\texttt{\textbackslash N}, \texttt{\textbackslash diff}, \texttt{\textbackslash e}, \texttt{\textbackslash lra}), with extra wrapper packages (bm, bbm, dsfont, xspace, cancel, mathtools, enumitem). The delivered numbering is the unit's own. Assets: no retained span contains a figure environment, a graphics-inclusion command, a PDF-page inclusion, a table or a TikZ picture; no image file is copied into this package, the package declares no asset dependency and no image file is redistributed with it. Licence: version arXiv:2408.15030v1 of this article is distributed under the Creative Commons Attribution 4.0 International licence, as recorded in the arXiv licence metadata for this exact version and shown on the abstract page https://arxiv.org/abs/2408.15030v1. A full rights notice is delivered at licenses/arxiv-2408.15030-CC-BY-4.0.txt. Characters: every retained excerpt is pure ASCII, so the delivered engine, XeLaTeX with its default Latin Modern text font, reports no missing character.
\begin{lemma}[Gr\"unbaum's inequality on intervals; $N<-1$, $N=\infty$]\label{lm:1d_neg}
Let $-\infty \le a<0<b \le \infty$ and $w\colon (a,b) \lra [0,\infty)$ be a nonnegative integrable function satisfying
$\int_a^b w(x)\diff x=1$ and $\int_a^b xw(x)\diff x=0$.
\begin{enumerate}[{\rm (i)}]
\item\label{1d_neg}
If $w^{1/(N-1)}$ is convex for some $N<-1$, then we have
\begin{equation}\label{eq:1d_neg}
\int_a^0 w(x) \diff x \geq \biggl( \frac{N}{N+1} \biggr)^N, \qquad
\int_0^b w(x) \diff x \geq \biggl( \frac{N}{N+1} \biggr)^N.
\end{equation}
\item\label{1d_inf}
If $\log w$ is concave, then we have
\begin{equation}\label{eq:1d_inf}
\int_a^0 w(x) \diff x \geq \e^{-1}, \qquad
\int_0^b w(x) \diff x \geq \e^{-1}.
\end{equation}
\end{enumerate}
\end{lemma}

\begin{equation}\label{eq:H_neg}
R(x)^{1/N} \ge R(0)^{1/N} \biggl( 1+\frac{c}{N}x \biggr), \qquad
c:=\frac{w(0)}{R(0)}>0.
\end{equation}

\begin{align}
b &= \int _a^b R(x)\diff x
 \le  \int_a^{1/c} R(0)\biggl( 1+\frac{c}{N}x \biggr)^N \diff x +b-\frac{1}{c} \label{eq:b_neg}\\
&= R(0) \Biggl[ \frac{N}{c(N+1)} \biggl( 1+\frac{c}{N}x \biggr)^{N+1} \Biggr]_a^{1/c} +b-\frac{1}{c} \nonumber\\
&\le \frac{R(0)}{c} \biggl(\frac{N+1}{N}\biggr)^N +b-\frac{1}{c}, \nonumber
\end{align}

\begin{lemma}[Rigidity on intervals; $N<-1$, $N=\infty$]\label{lm:=_inf}
\begin{enumerate}[{\rm (i)}]
\item\label{1d=_neg}
If equality holds in the former inequality in \eqref{eq:1d_neg},
then, for some $c>0$, we have $a=-\infty$, $b=1/c$, and
\begin{equation}\label{eq:=_inf}
w(x) =c \biggl( \frac{N}{N+1} \biggr)^N \biggl( 1+\frac{c}{N}x \biggr)^{N-1}.
\end{equation}
Similarly, if equality holds in the latter inequality in \eqref{eq:1d_neg},
then, for some $c>0$, we have $a=-1/c$, $b=\infty$, and
\[
w(x) =c \biggl( \frac{N}{N+1} \biggr)^N \biggl( 1-\frac{c}{N}x \biggr)^{N-1}.
\]

\item\label{1d=_inf}
If equality holds in the former inequality in \eqref{eq:1d_inf},
then we have $a=-\infty$, $b=1/c$, and $w(x)=c \e^{cx-1}$ for some $c>0$.
Similarly, if equality holds in the latter inequality in \eqref{eq:1d_inf},
then we have $a=-1/c$, $b=\infty$, and $w(x)=c \e^{-cx-1}$ for some $c>0$.
\end{enumerate}
\end{lemma}

\begin{proof}
We will consider only the former inequalities in \eqref{eq:1d_neg} and \eqref{eq:1d_inf},
since the latter inequalities can be handled by reversing the interval.

\eqref{1d=_neg}
We first assume $b<\infty$.
In this case, we deduce from the proof of Lemma~\ref{lm:1d_neg} that $a=-\infty$, $b=1/c$, and that equality holds in \eqref{eq:H_neg}.
Combining these with $R(1/c)=1$ implies
\[ R(x) =\biggl( \frac{N}{N+1} \biggr)^N \biggl( 1+\frac{c}{N}x \biggr)^N, \]
and hence
\[ w(x) =R'(x) =c \biggl( \frac{N}{N+1} \biggr)^N \biggl( 1+\frac{c}{N}x \biggr)^{N-1}. \]

Next, we show that equality never holds when $b=\infty$.
Assume in contrary that equality holds in the former inequality in \eqref{eq:1d_neg}.
For (large) $k \in \N$, we take $a_k \in (a,0)$ and $w_k$ as in the proof of Lemma~\ref{lm:1d_neg}, and put
\[ R_k(x) :=\int_{a_k}^x w_k(s) \diff s, \qquad
 c_k:=\frac{R'_k(0)}{R_k(0)} =\frac{w_k(0)}{\int_{a_k}^0 w_k(s) \diff s}. \]
Observe from
\[ \lim_{k \to \infty} c_k = \frac{w(0)}{\int_a^0 w(s) \diff s} =:c \]
that $k \ge 1/c_k$ for large $k$.
Then, in the estimate \eqref{eq:b_neg} in the proof of Lemma~\ref{lm:1d_neg},
\[ k-\frac{1}{c_k} -\int_{1/c_k}^k R_k(x) \diff x =\int_{1/c_k}^k \bigl( 1-R_k(x) \bigr) \diff x \]
necessarily tends to $0$ as $k \to \infty$.
This implies that $\int_a^{1/c} w(s) \diff s =1$,
which is a contradiction since we assumed $w>0$ on $(a,b)$.

\eqref{1d=_inf}
Let us begin with a direct proof of \eqref{eq:1d_inf} under $b<\infty$.
We use the same notations as Lemma~\ref{lm:1d_neg}.
It follows from the concavity of $\log w$ that
\begin{align*}
R\bigl( (1-\lambda)x +\lambda y \bigr)
&= \int_0^1 w\bigl( \theta(t) \bigr) \theta'(t) \diff t \\
&\ge \int_0^1 w\bigl( \sigma(t) \bigr)^{1-\lambda} w\bigl( \tau(t) \bigr)^{\lambda}
 \biggl( (1-\lambda) \frac{R(x)}{w(\sigma(t))} +\lambda \frac{R(y)}{w(\tau(t))} \biggr) \diff t \\
&= \int_0^1 \biggl\{ (1-\lambda) R(x) \biggl( \frac{w(\tau(t))}{w(\sigma(t))} \biggr)^{\lambda}
 +\lambda R(y) \biggl( \frac{w(\sigma(t))}{w(\tau(t))} \biggr)^{1-\lambda} \biggr\} \diff t.
\end{align*}
Together with the concavity of $\log$ and Jensen's inequality, we obtain the concavity of $\log R$:
\begin{align*}
\log R\bigl( (1-\lambda)x +\lambda y \bigr)
&\ge (1-\lambda) \log\biggl[ R(x) \int_0^1 \biggl( \frac{w(\tau(t))}{w(\sigma(t))} \biggr)^{\lambda} \diff t \biggr] \\
&\quad +\lambda \log\biggl[ R(y) \int_0^1 \biggl( \frac{w(\sigma(t))}{w(\tau(t))} \biggr)^{1-\lambda} \diff t \biggr] \\
&\ge (1-\lambda) \biggl( \log R(x) +\int_0^1 \lambda \log\biggl[ \frac{w(\tau(t))}{w(\sigma(t))} \biggr] \diff t \biggr) \\
&\quad +\lambda \biggl( \log R(y) +\int_0^1 (1-\lambda) \log\biggl[ \frac{w(\sigma(t))}{w(\tau(t))} \biggr] \diff t \biggr) \\
&= (1-\lambda) \log R(x) +\lambda \log R(y).
\end{align*}
Therefore,
\[ \log R(x) \le \log R(0) +cx, \qquad c:=\frac{w(0)}{R(0)}>0. \]
When $b \ge 1/c$, we have
\[
b \le \int_a^{1/c} R(0) \e^{cx} \diff x +b-\frac{1}{c}
 \le \frac{R(0)}{c} \e +b-\frac{1}{c},
\]
which yields \eqref{eq:1d_inf}.
In the case of $b<1/c$, since $1 \le R(0) \e^{cb} \le R(0) \e^{cx}$ for $x \ge b$,
we similarly find
\[
b \le \int_a^{1/c} R(0) \e^{cx} \diff x -\biggl( \frac{1}{c}-b \biggr)
 \le \frac{R(0)}{c} \e +b-\frac{1}{c}.
\]

If equality holds, then we have $a=-\infty$, $b=1/c$, and $\log w$ is an affine function.
Put $w(x) =\alpha \e^{\delta x}$ ($\alpha,\delta>0$) and observe that
\[ R(b) =\frac{\alpha}{\delta}\e^{\delta/c} =1, \qquad
 c =\frac{w(0)}{R(0)} =\delta. \]
Therefore, we obtain $\delta=c$, $\alpha=c\e^{-1}$, and $w(x)=c\e^{cx-1}$.
One can also show that equality never holds when $b=\infty$ in a similar way to \eqref{1d=_neg}.
\end{proof}

\end{document}
